% Русская версия. Структура секций, формул и ссылок должна быть идентична
% main.tex.
\documentclass[10pt,twocolumn]{article}
\usepackage[T2A]{fontenc}
\usepackage[utf8]{inputenc}
\usepackage[russian,english]{babel}
\usepackage[margin=0.72in]{geometry}
\usepackage{amsmath,amssymb,booktabs,graphicx,microtype,url}
\usepackage{tikz}
\usetikzlibrary{arrows.meta,positioning,fit}
\microtypesetup{expansion=false}
\usepackage[hidelinks]{hyperref}

\title{Один глобальный луч на множестве GPU: Точный высокопроизводительный beam search с миллиардным frontier}
\author{Иван Литвак\\Независимый исследователь\\\href{mailto:vanya.litvak@gmail.com}{vanya.litvak@gmail.com}}
\date{}

\newcommand{\beam}{B}
\newcommand{\frontier}{\mathcal{F}}
\newcommand{\cand}{\mathcal{C}}

\begin{document}
\selectlanguage{russian}
\maketitle

\begin{abstract}
Мы представляем первую опубликованную высокопроизводительную multi-GPU систему
beam search, которая физически шардирует один логический frontier между
устройствами и совместно строит один следующий луч. Генерация кандидатов,
оценка, дедупликация по 128-битному ключу, маршрутизация и хранение frontier
остаются на GPU; на host поступают только компактные управляющие данные и
метаданные родословной. При отсутствии семантических локальных ограничений
шардов потоковая процедура возвращает тот же глобальный top-$B$ по
редуцированной оценке ключа, что и монолитная реализация, при той же квантизации
оценки, порядке разрешения равных оценок для фиксированной раскладки и
заявленной 128-битной модели идентичности. На восьми H200 запуск Cube4
завершил одну насыщенную глубину с физически распределённым удерживаемым
frontier из $B_{\rm eff}=2{,}900{,}361{,}216$ состояний за $931.266$ с.
При 24 генераторах это $69{,}608{,}669{,}184$ номинальные пары
родитель--действие до дедупликации, или $74.746$ миллиона пар/с end to end.
Запуск был остановлен во время следующей глубины и не решил задачу. Отдельно
измеренный запуск на двух T4 при $B_{\rm eff}=82{,}837{,}504$ имеет
$1{,}988{,}100{,}096$ номинальных потомков на насыщенную глубину; глубины
7--10 занимали в среднем $65.670$ с, то есть $30.274$ миллиона логических
потомков/с end to end. Задачи и модели в этих запусках различаются, поэтому
они не устанавливают strong scaling. Новизна состоит в
полной системной совокупности свойств, а не в параллельном, GPU- или
распределённом beam search по отдельности.
\end{abstract}

\section{Введение и постановка задачи}
Beam search хранит ограниченный набор текущих состояний --- луч. На каждом
шаге он применяет допустимые ходы, объединяет повторяющиеся состояния и
оставляет не более $B$ записей с наименьшими оценками. Если различных ключей
меньше $B$, сохраняются все. Владелец --- GPU, назначенная по ключу состояния:
у неё встречаются копии с разных устройств. Точный top-$B$ означает совпадение
при тех же квантованных оценках, 128-битных ключах и порядке равных оценок для
фиксированной раскладки; он не гарантирует кратчайшего решения, поскольку
beam search отбрасывает часть направлений.
На рис.~\ref{fig:beam_step} небольшой пример проходит объединение повторов,
глобальный отбор и отдельный случай равных оценок на границе.

\begin{figure}[t]
\centering
\setlength{\fboxsep}{4pt}
\begin{tabular}{c@{\hspace{4pt}}c}
\fbox{\begin{tabular}{c}GPU 0\\$A{:}2,\ B{:}7,\ C{:}4$\end{tabular}} &
\fbox{\begin{tabular}{c}GPU 1\\$A{:}2,\ D{:}3,\ E{:}8$\end{tabular}}\\[5pt]
\multicolumn{2}{c}{$\Downarrow$ один владелец для $A$}\\[3pt]
\multicolumn{2}{c}{\fbox{$A{:}2,\ B{:}7,\ C{:}4,\ D{:}3,\ E{:}8$}}\\[5pt]
\multicolumn{2}{c}{$\Downarrow\quad B=3\quad\Downarrow$}\\[3pt]
\multicolumn{2}{c}{\fbox{$A{:}2,\ D{:}3,\ C{:}4$}}\\[7pt]
\multicolumn{2}{c}{\begin{tabular}{c}
\scriptsize Пример границы: $1,2,3,3$ при $B=3$ и $T=3$\\
\scriptsize обе $3$ сохраняются при $s\leq T$\\
\scriptsize одно место выбирает правило равенства
\end{tabular}}
\end{tabular}
\caption{Один условный шаг beam search при $B=3$: меньшая оценка лучше.
Одинаковые ключи объединяются между GPU до глобального отбора. Две записи
$A$ занимают одно место; идентичность здесь задаётся моделью 128-битных ключей.
На границе с равными оценками обе записи переживают отсечение до финального отбора.}
\label{fig:beam_step}
\end{figure}

Beam search становится задачей распределённых систем, когда удерживаемый
frontier перестаёт помещаться в память одной GPU. Расширение, оценка,
дедупликация, глобальный отбор и родословная должны распределяться без
превращения одного поиска в независимые реплики и без переноса frontier через
host memory. Мы исследуем, может ли GPU-кластер сохранить семантику одного
монолитного луча, масштабируя его ёмкость суммарной памятью устройств.

Пусть $G=(V,E)$ --- неявный ориентированный граф, $A_a(x)$ --- потомок
состояния $x$ под действием генератора $a\in\{1,\ldots,m\}$, а $\sigma(x,a)$
--- детерминированная вещественная функция оценки. Математический контракт и
стримы~2--5 допускают обученную или заданную вручную функцию, если стрим~1
выдаёт детерминированные квантованные ключи; поставляемые production-backend'ы
стрима~1 и представленные эксперименты являются нейросетевыми. Реализация
ранжирует только отсечённый и квантованный целочисленный ключ $s(x,a)$,
определённый в
формуле~\eqref{eq:score}. Скалярная оценка состояния является частным случаем
$\sigma(x,a)=q(A_a(x))$, где $q:V\rightarrow\mathbb{R}$, тогда как оценщик с
$m$ выходами выдаёт все оценки действий из одного вычисления родителя.
На глубине $d$ beam search вычисляет
\begin{align}
  \cand_d &= \{(x,a,A_a(x)):x\in\frontier_{d-1},\ a\in[m]\},\nonumber\\
  U_d &= \operatorname{Reduce}_{z}
  (\cand_d;\,\operatorname{lexmin}(s,\mathrm{payload})),\\
  \frontier_d &= \operatorname{Top}_{\beam}
  \left(U_d;\,s_d^*,\tau_L\right),
  \label{eq:beam}
\end{align}
где $z$ --- 128-битный ключ дочернего состояния,
$s_d^*(z)=\min\{s(x,a):z(A_a(x))=z\}$, а $\tau_L$ задаёт
детерминированный порядок разрешения равных оценок, индуцированный фиксированной раскладкой
рангов и шардов $L$. Таким образом, до глобального среза дедупликация сохраняет
лучшего кандидата для каждого ключа. Для насыщенного слоя, в котором к каждому удерживаемому
родителю применяются все $m$ генераторов, номинальное число кандидатов до
дедупликации равно $m\beam_{\rm eff}$. Поскольку runtime-выравнивание даёт
$\beam_{\rm eff}\geq\beam_{\rm req}$, измеренный H200-frontier
$\beam_{\rm eff}=2{,}900{,}361{,}216$ при $m=24$ означает
$69{,}608{,}669{,}184$ номинальных логических потомка на одной насыщенной
глубине. Материализация полных логических состояний-потомков Cube4 по 96 байт
потребовала бы более $6.68$ ТБ ещё до метаданных; система избегает такого переноса на host
и глобальной сортировки на одном устройстве.

Наше требование сильнее, чем GPU-ускоренное вычисление эвристики. В течение
одного уровня массивы кандидатов и frontier из формулы~\eqref{eq:beam} остаются
в памяти GPU. CPU получает только $O(p)$ управляющих счётчиков и компактные
метаданные отобранных кандидатов, но не полные массивы состояний или frontier.
Термин ``один согласованный поиск'' имеет операционное значение: каждая GPU
участвует в формуле~\eqref{eq:beam}, а ранги совместно фиксируют непересекающиеся
шарды одного глобально определённого следующего frontier. Это не ансамбль
независимых лучей или реплик.

На рис.~\ref{fig:global_beam} показано, где остаются полные данные кандидатов
и где на CPU поступают только компактные управляющие метаданные. Отдельной
квоты выживших для каждой GPU нет: все ранги участвуют в общем срезе.

\begin{figure*}[t]
\centering
\begin{tikzpicture}[>=Stealth, every node/.style={font=\small,align=center},
 box/.style={draw,rounded corners,text width=2.05cm,minimum height=1.05cm,font=\scriptsize},
 node distance=0.3cm]
\node[box] (front) {$\frontier_{d-1}$\\один логический луч};
\node[box,right=of front] (batch) {GPU 0, 1, 2, $\ldots$\\локальные батчи};
\node[box,right=of batch] (route) {компактные записи\\владелец $o(z)$};
\node[box,right=of route] (reduce) {редукция владельца\\общий порог $T^*$};
\node[box,right=of reduce] (next) {$\frontier_d$\\новый общий луч};
\draw[->] (front) -- (batch);
\draw[->] (batch) -- (route);
\draw[->] (route) -- (reduce);
\draw[->] (reduce) -- (next);
\node[draw,dashed,below=0.65cm of reduce,font=\scriptsize,align=center] (cpu)
 {CPU: компактные счётчики и метаданные родословной};
\draw[->,dashed] (reduce) -- (cpu);
\node[draw,dotted,fit=(front)(batch)(route)(reduce)(next),inner sep=0.12cm] {};
\end{tikzpicture}
\caption{Один глобальный луч физически распределён по GPU. Одинаковые ключи
поступают одному владельцу; порог и финальный top-$B$ относятся ко всему
множеству кандидатов, без семантических локальных квот. Сплошные стрелки
обозначают GPU hot data, пунктирная --- только компактные метаданные на CPU.}
\label{fig:global_beam}
\end{figure*}

\paragraph{Вклад.}
Система даёт (i) GPU-резидентный hot-data контур, распределяющий один логический
луч без переноса frontier через host memory; (ii) точную семантику монолитного
глобального top-$B$ по редуцированной оценке ключа при той же квантизации
оценки, порядке разрешения равных оценок для фиксированной раскладки и
заявленной 128-битной модели идентичности, без семантических локальных
ограничений шардов; и (iii) измеренную насыщенную глубину с миллиардным
frontier на восьми H200 и отдельный production-результат на двух T4.
В совокупности это первая опубликованная
высокопроизводительная multi-GPU система beam search с одним физически
шардированным глобальным frontier и заявленным контрактом точности. Реализация
допускает не более 128 рангов; исполнение за пределами восьми GPU и strong
scaling на фиксированной нагрузке пока не измерены.

\section{Предшествующие работы и пробел}
Параллельный beam search появился задолго до GPU: специализированные
speech-ускорители исследовались уже в 1986 году
~\cite{anantharaman1986accelerator}. Затем Wijs и Lisser распределили один
глобальный луч между CPU~\cite{wijs2007distributed}, а Chong и соавторы
разместили data-parallel speech beam на одной GPU~\cite{chong2008lvcsr}.
Следовательно, сами по себе параллельный, распределённый и GPU beam search не
новы. Современная CPU-работа удерживает лучи до 20 миллионов на общей памяти, а
MPI-workers запускают независимые рандомизированные лучи
~\cite{frohner2022parallel}; Kuroiwa и Beck также исследуют глобальные и
локально разделённые CPU-варианты~\cite{kuroiwa2024parallel}.

В GPU- и эвристическом комбинаторном поиске GA* задаёт single-GPU A* без нейросетевого
оценщика~\cite{zhou2015gpuastar}, GPUAStar удерживает neural A* на одной H200
~\cite{soltani2026gpuastar}, а HDA* задаёт hash-distributed CPU best-first
search~\cite{kishimoto2013hda}. Ближайшим алгоритмическим предшественником
является CayleyPy: его агенты для больших кубов используют одну GPU и
$B=2^{24}$~\cite{chervov2025large,soibelman2025cayleypy}. AlphaCube и
DeepCubeA являются соответственно однокарточным beam search и пакетным
weighted A*~\cite{takano2023alphacube,agostinelli2019deepcubea}.

Ближайший смежный системный результат --- QiankunNet-cuSCI: он генерирует
конфигурации, выполняет распределённую глобальную дедупликацию, применяет
нейросетевой Transformer и выбирает top-$K$ на 64 A100 с заявленной
эффективностью strong scaling 91.06\%~\cite{sun2026cunnqssci}. Поэтому нельзя
заявлять приоритет общего конвейера expand--deduplicate--infer--select, но работа
не представлена как beam search. Она удерживает $K=32{,}000$ или $256{,}000$
отобранных SCI-конфигураций примерно из миллиарда сырых кандидатов, а её
глобальные барьеры переносят полные данные кандидатов через host memory.

Другие результаты миллиардного масштаба измеряют иные объекты. GPUexplore 3.0
обходит более 3.15 миллиарда достижимых состояний на четырёх GPU, но выполняет
полный обход без ограниченного top-$B$ с ранжированием по оценке~\cite{wijs2024gpuexplore}. PathWeaver
конвейерно обрабатывает приближённые запросы ближайших соседей по сохранённому
графу~\cite{kim2025pathweaver}; multi-GPU speech- и языковые декодеры
распределяют независимые высказывания, предложения или исполнение модели вокруг
малых лучей, но физически не шардируют один frontier состояний
~\cite{junczysdowmunt2016marian,braun2020gpuasr,chan2025trie}.

Таблица~\ref{tab:prior_art} применяет единую таксономию масштаба: удерживаемая
ширина луча не равна сырым кандидатам, накопленным расширениям, размеру графа
или пакету независимых запросов. Предшествующие системы устанавливают важные
подмножества свойств параллельного beam search, исполнения на GPU,
распределённого поиска или крупномасштабного исследования состояний, но не
сообщают об их полной совокупности в одной системе. Эта работа представляет
первую опубликованную высокопроизводительную multi-GPU систему beam search,
которая объединяет один физически шардированный логический frontier,
GPU-резидентный hot-data контур кандидатов и frontier, глобальную дедупликацию
по 128-битному ключу, отсутствие семантических локальных ограничений шардов и
точный монолитный глобальный top-$B$ по редуцированной оценке ключа при
заявленных квантизации, порядке разрешения равных оценок и модели идентичности
при измеренном миллиардном удерживаемом frontier. Заявление о приоритете
относится только к этой полной системной совокупности.

\begin{table*}[t]
\centering
\footnotesize
\caption{Ближайшие предшествующие системы в единой таксономии масштаба.
``Иной масштаб'' намеренно не интерпретируется как удерживаемая ширина луча.}
\label{tab:prior_art}
\setlength{\tabcolsep}{2.5pt}
\begin{tabular}{@{}p{2.7cm}p{3.7cm}p{3.0cm}p{3.0cm}p{4.2cm}@{}}
\toprule
Система & Платформа и объект поиска & Удерживаемый луч & Иной масштаб & Чего нет из заявленной совокупности \\
\midrule
Wijs--Lisser 2007 & distributed CPU; один глобальный луч & 50,000 & 135.96M накопленных состояний & GPU-контур и device-resident scorer \\
Chong et al. 2008 & одна 8800 GTX; активное speech-множество & предел 20,000 & модель на 1M состояний & multi-GPU frontier состояний \\
Frohner et al. 2022 & 46 CPU threads; один shared-memory луч & 20,000,000 & MPI: независимые лучи & GPU и распределённое исполнение одного луча \\
CayleyPy 2025 & одна GPU; один луч графа Кэли & $2^{24}$ & независимые агенты & multi-GPU глобальный frontier \\
GPUexplore 3.0 & четыре V100; полное множество достижимых & не луч & $>3.15$B состояний & ранжированная оценка и top-$B$ \\
PathWeaver 2025 & четыре A6000; независимые ANNS-запросы & не сообщён & датасет 50M векторов & точная послойная семантика и большой луч \\
QiankunNet-cuSCI & 64 A100; отобранные SCI-конфигурации & не луч; $K=32{,}000/256{,}000$ & около 1B сырых кандидатов & не заявлено как beam search; существенно меньший удерживаемый $K$; перенос полных кандидатов через host \\
Эта работа & восемь H200; один шардированный глобальный луч & 2,900,361,216 effective & 69.609B номинальных сырых пар/глубину & глубина 8 завершена; полное решение и strong scaling на одной нагрузке не измерены \\
\bottomrule
\end{tabular}
\end{table*}

\section{Один уровень поиска в пяти GPU-стримах}
Пять стримов перекрывают работу над разными микробатчами, а не образуют пять
последовательных проходов. Рис.~\ref{fig:overlap} --- концептуальная схема,
не трасса профилировщика и не измерение длительности отдельных стадий.

\begin{figure*}[t]
\centering
\begin{tikzpicture}[x=1.05cm,y=0.48cm,>=Stealth,
 every node/.style={font=\scriptsize}]
\foreach \name/\y in {$S_1$/5,$S_2$/4,$S_3$/3,$S_4$/2,$S_5$/1}
  {\node[anchor=east] at (0,\y) {\name}; \draw[gray!50] (0.2,\y) -- (8.2,\y);}
\draw[fill=gray!25] (0.3,4.7) rectangle (2.2,5.25);
\draw[fill=gray!25] (2.4,4.7) rectangle (4.25,5.25);
\draw[fill=gray!25] (0.3,3.7) rectangle (2.0,4.25);
\draw[fill=gray!25] (2.4,3.7) rectangle (4.1,4.25);
\draw[fill=gray!45] (2.45,2.7) rectangle (3.55,3.25);
\draw[fill=gray!45] (4.5,2.7) rectangle (5.6,3.25);
\draw[fill=gray!65] (3.7,1.7) rectangle (5.0,2.25);
\draw[fill=gray!65] (5.85,1.7) rectangle (7.1,2.25);
\draw[fill=gray!65] (3.7,0.7) rectangle (5.0,1.25);
\draw[fill=gray!65] (5.85,0.7) rectangle (7.1,1.25);
\foreach \x/\y/\tag in {1.25/5/$b_1$,3.3/5/$b_2$,1.15/4/$b_1$,3.25/4/$b_2$,3.0/3/$b_1$,5.05/3/$b_2$,4.35/2/$b_1$,6.48/2/$b_2$,4.35/1/$b_1$,6.48/1/$b_2$}
  {\node at (\x,\y) {\tag};}
\draw[->] (2.2,4.6) -- (2.45,3.3);
\draw[->] (2.0,3.6) -- (2.45,3.3);
\draw[dashed] (7.35,0.55) -- (7.35,5.45);
\node[anchor=west] at (7.43,3) {drain};
\end{tikzpicture}
\caption{Концептуальная схема перекрытия, не измеренная трасса. Стримы 1 и 2
оценивают батч и вычисляют его ключи до потребления стримом 3; следующие батчи
обрабатываются, пока предыдущие редуцируются или передаются. Стрим 5
обменивается записями с владельцами; коллектор стрима 4 обрабатывает локальные
и полученные записи. Финальный отбор следует после drain. Длины полос не
кодируют измеренное время.}
\label{fig:overlap}
\end{figure*}

На входе находится глобальный луч $\frontier_{d-1}$, распределённый между $p$
GPU. Если GPU $r$ владеет $N_r$ родителями, а граф имеет $m$ ходов, она
порождает $mN_r$ логических кандидатов. Для $y=A_a(x)$ до материализации
кандидат представлен как
\begin{equation}
 c=(z(y),s(x,a),i_{\mathrm{parent}},a,r_{\mathrm{src}}),
 \label{eq:candidate}
\end{equation}
где $z(y)\in\{0,1\}^{128}$ --- хеш дочернего состояния,
$s(x,a)$ --- квантованная оценка кандидата, а
$r_{\mathrm{src}}$ указывает GPU, хранящую родителя. Временное дочернее
состояние может создаваться для проверки цели или скалярной оценки; полный
расширенный слой потомков не хранится. Состояния следующего луча
восстанавливаются после глобального отбора.

\paragraph{Стримы 1--2: оценка и идентичность.}
Родители делятся на микробатчи и помещаются в переиспользуемые слоты кольцевого
буфера. Стримы~1 и~2 одновременно потребляют один слот. Стрим~1 вычисляет
выбранную функцию оценки и записывает только
\begin{equation}
 s(x,a)=\operatorname{rint}_{\mathrm{RN-even}}
 \left(S\,\operatorname{clip}(\sigma(x,a),0,q_{\max})\right)
 \in\mathbb{N},
 \label{eq:score}
\end{equation}
но не её вещественный выход. Здесь RN-even --- используемое GPU-контуром
округление к ближайшему с разрешением половин к чётному. Скалярный путь оценивает дочерние состояния, а
путь с $m$ выходами получает все $m$ оценок действий для каждого родителя.
Стрим~2 применяет каждый ход в регистрах, проверяет цель и записывает 128-битный
хеш Зобриста. Поэтому вычисление оценки и преобразование состояний перекрываются без
хранения $mN_r$ дочерних состояний.

\paragraph{Стрим 3: раннее сокращение и владение.}
Когда группа кольцевых слотов имеет готовые оценки и хеши, стрим~3 удаляет
оценки выше текущего порога $T$, компактизирует оставшиеся записи, сортирует их
по $z$ и объединяет равные хеши. Для каждого хеша сохраняется лексикографически
минимальная пара $(s,\mathrm{payload})$, где payload восстанавливает исходного
родителя и ход. Затем уникальному кандидату назначается владелец
\begin{equation}
 o(x)=\operatorname{mix}(z(x))\bmod p.
 \label{eq:owner}
\end{equation}
Локальные кандидаты поступают в локальный коллектор, остальные группируются по
владельцу. Это первое сокращение перед коммуникацией уменьшает трафик, но не
заменяет последующую глобальную дедупликацию.

\paragraph{Стрим 5: обмен между GPU.}
Стрим~5 напрямую обменивает сгруппированные по владельцу метаданные между GPU
через NCCL. Правило владения даёт
\begin{equation}
 z(x)=z(y)\Longrightarrow o(x)=o(y),
 \label{eq:owner_invariant}
\end{equation}
поэтому дубликаты с разных исходных GPU в итоге встречаются у одного владельца.
Стрим~5 выполняет только коммуникацию; полученные кандидаты поступают в тот же
коллектор, что и локальные.

\paragraph{Стрим 4: постоянное множество выживших.}
Коллектор распределяет принадлежащих GPU кандидатов по логическим шардам.
Каждый шард имеет два постоянных буфера: стрим~3 дописывает в один, пока
стрим~4 очищает другой. Задача стрима~4 объединяет чистую и недавно дописанную
области, применяет снимок $T$, компактизирует, сортирует по $z$ и сохраняет
лучшие метаданные для каждого хеша. Внутри шарда нет top-$k$ и семантической
квоты: шард может сохранить всех кандидатов с $s\leq T$. Поэтому шардирование не
аппроксимирует срез по оценке; только порядок границы равных оценок задаётся
фиксированной раскладкой.

\paragraph{Классический top-$\beam$ и монотонное пороговое отсечение.}
Пусть $C$ --- полное мультимножество кандидатов, а $s(c)$ --- квантованная
целочисленная оценка из формулы~\eqref{eq:score}. Тогда
$U=\operatorname{Reduce}_{z}(C;\operatorname{lexmin}(s,\mathrm{payload}))$ ---
его редукция до одной записи на ключ,
которая несёт $s^*(z)=\min_{c:z(c)=z}s(c)$ и соответствующий лучший payload.
Для фиксированной раскладки $L$ ключи упорядочиваются по
$\kappa_L(z)=(s^*(z),\tau_L(z))$, где меньшая оценка лучше, а $\tau_L$ ---
детерминированное правило реализации: сначала ранг, затем физический порядок
для равных оценок. Классическая эталонная реализация материализует $U$,
сортирует его по $\kappa_L$ и при $b=\min(\beam,|U|)$ возвращает
\begin{equation}
 \operatorname{Top}_{\beam}(U)=
 \{z\in U:\kappa_L(z)\leq \kappa_{L,(b)}(U)\}.
 \label{eq:classical_topb}
\end{equation}
Наша реализация никогда не материализует $U$. После обновления $j$ пусть
$C_j\subseteq C_{j+1}\subseteq C$ --- уже обработанные кандидаты, а
$U_j=\operatorname{Reduce}_{z}(C_j)$. Область ключей $U_j$ может только расти,
а оценка существующего ключа $s_j^*(z)$ --- только уменьшаться. Как только
$|\operatorname{dom}U_j|\geq\beam$, гистограммы текущих редуцированных оценок
задают
\begin{align}
 \widehat T_j&=\min\left\{t:\sum_{u=0}^{t}H_j(u)\geq\beam\right\},\nonumber\\
 T_j&=\min\{T_{j-1},\widehat T_j\},\qquad T_0=+\infty.
 \label{eq:online_threshold}
\end{align}
Пока $|\operatorname{dom}U_j|<\beam$, полагается $\widehat T_j=+\infty$.
Стримы~3 и~4 могут удалять лишь кандидатов с $s(c)>T_j$; граница
$s(c)=T_j$ сохраняется, поскольку окончательное попадание в луч на ней зависит
от $\tau_L$.

\paragraph{Точность.}
Обработка следующего кандидата либо вводит новый ключ, либо уменьшает лучшую
оценку существующего, поэтому $\beam$-я наименьшая редуцированная оценка не может
увеличиться; следовательно, $\widehat T_{j+1}\leq\widehat T_j$ и
$T_{j+1}\leq T_j$. Если кандидат $c$ удалён при обновлении $j$, то
$s(c)>T_j$, а $U_j$ уже содержит не менее $\beam$ различных ключей с текущими
лучшими оценками не выше $T_j$. Эти ключи останутся различными, а их финальные
лучшие оценки могут только уменьшиться. Если $z(c)$ входит в их число, то $c$
хуже уже сохранённого представителя; иначе не менее $\beam$ других ключей
строго лучше. Поэтому $c$ не может определять отобранный финальный минимум.
Обратно, кандидат, достигающий финального минимума ключа из глобального
top-$\beam$, не может быть удалён: для этого уже должны существовать $\beam$
различных ключей со строго меньшими оценками. Следовательно, пороговая схема
сохраняет точно тот же классический top-$\beam$ по редуцированным оценкам ключей
при той же квантизации для фиксированной $L$ с тем же правилом разрешения
равных оценок на границе.

\paragraph{Модель идентичности и область точности.}
Доказательство является точным для квантизации оценки из
формулы~\eqref{eq:score}, множества, индексированного 128-битным идентификатором
Зобриста, и полного фиксированного порядка $\kappa_L$.
Контур кандидатов не
переносит все 120 логических байт состояния, поэтому равные 128-битные ключи
считаются равными состояниями без полного сравнения состояний. Следовательно,
отбор математически точен в явно заданной модели ключа; остаточный риск связан
с коллизией 128-битного хеша, а не с аппроксимацией при шардировании,
пороговом отсечении или top-$\beam$. Изменение числа GPU или шардов может
изменить $\tau_L$ и тем самым набор состояний с равной граничной оценкой;
текущая реализация не заявляет инвариантность выбранных ID к разбиению.

\paragraph{Финальный срез и материализация.}
После осушения всех пяти стримов ранги объединяют итоговые гистограммы. Если
$|U|\geq\beam$, они выбирают минимальный $T^*$, для которого
\begin{equation}
 \sum_{u=0}^{T^*}H(u)\geq\beam.
 \label{eq:threshold}
\end{equation}
Все состояния ниже $T^*$ сохраняются, а детерминированное правило обрезает его
границу ровно до $\beam$ состояний. Если $|U|<\beam$, вместо этого сохраняются
все доступные ключи. Поэтому следующий frontier всегда содержит ровно
$\min(\beam,|U|)$ состояний. Отобранные метаданные равномерно распределяются
между GPU. Только после этого каждая исходная GPU восстанавливает
запрошенные дочерние состояния и отправляет их непосредственно целевым GPU. CPU
получает $O(p)$ векторов счётчиков и компактные метаданные отобранных кандидатов
для восстановления родословной, но не полные массивы состояний или frontier.

Пять стримов управляются готовностью данных, а не выполняются как пять
последовательных фаз. Стримы~1--2 почти непрерывно заполняют кольцевые слоты,
пока стримы~3 и~4 сокращают предыдущие батчи, а стрим~5 обменивает удалённых
выживших. Математический инвариант состоит в том, что каждый хеш, переживший
финальный срез, объединён у своего единственного владельца до вычисления
формулы~\eqref{eq:threshold}. Поэтому для фиксированной раскладки распределённое
выполнение реализует тот же глобальный top-$\beam$ по редуцированной оценке
ключа, что и монолитный эталон при той же квантизации, не материализуя сырой
слой кандидатов.

\section{Эксперименты}
\paragraph{Протокол.}
Мы оцениваем вычислительную стоимость, а не качество решения, и приводим
измерение только вместе с оборудованием, нагрузкой, контрактом модели и
областью. Production-время включает генерацию, дедупликацию, коммуникацию,
отбор и историю; оно не взаимозаменяемо с изолированной скоростью стрима~1.
Пустая ячейка означает ``не измерено'', а не расчётную оценку. В
таблице~\ref{tab:platforms} собраны сильнейшие локально проверенные измерения.

\begin{table*}[t]
\centering
\footnotesize
\caption{Измеренная производительность по платформам. E2E --- критический путь
production-запуска, micro --- изолированный конвейер ядер, а ёмкость не содержит
заявления о времени. н/в: не вычислено; н/п: неприменимо; н/з: не записано.}
\label{tab:platforms}
\setlength{\tabcolsep}{3pt}
\begin{tabular}{@{}p{2.0cm}p{2.4cm}p{3.8cm}p{3.2cm}p{1.8cm}p{1.8cm}@{}}
\toprule
Платформа & Нагрузка & Луч $B$ & Измерение & с/глубину & Сырая скорость \\
\midrule
RTX 3070 8GB & задача 0, глубина 20 & $2^{22}$ & E2E завершён; всего $60.623$ с & 3.031 сред. & н/в \\
RTX 3070 8GB & стрим 1, $B_\mu=4096$, $c=3$ & --- & micro завершён & н/п & $48.142$ M/с \\
Kaggle $2\times$T4 & задача 20, радиус 4, предел 72 & $B_{\rm req}=82{,}615{,}524$; $B_{\rm eff}=82{,}837{,}504$ & решено; $350.402$ с; CSV+сводка & 65.670 насыщ. & $30.274$ M/с \\
$8\times$A100 40GB & задачи 900/950/1000, глубина 8 & $B_{\rm req}=770{,}883{,}178$; $B_{\rm eff}$ не записан & стабилен; $39{,}745$ MiB/GPU; сводка & н/з & н/з \\
$8\times$H200 & Cube4, задача 1000, глубина 8 & $B_{\rm req}=2.900$ млрд; $B_{\rm eff}=2{,}900{,}361{,}216$ & E2E-глубина завершена; $\sim141{,}820$ MiB/GPU занято & 931.266 & $74.746$ M/с \\
\bottomrule
\end{tabular}
\end{table*}

В таблице~\ref{tab:platforms} значение RTX равно полному E2E wall time,
делённому на 20 завершённых глубин, а значение Kaggle --- среднему для
насыщенных глубин 7--10 из сводки. T4 CSV фиксирует запрошенную и выровненную
эффективную ширину; скорость $30.274$ M/с явно вычислена как
$24\times82{,}837{,}504/65.670$ и считает логических потомков до дедупликации.
Для RTX E2E-скорость не выводится, поскольку ранние глубины не насыщены.
Микробенчмарки ядер и capacity-only сводка A100 не задают E2E-время глубины,
поэтому помечены как неприменимые или не записанные. H200-скорость --- это
номинальное $24B_{\rm eff}/931.266$; она считает пары родитель--действие, а не
различные состояния после дедупликации.

\paragraph{Локальная RTX 3070.}
На GPU ноутбука folded-input half2 ядро повысило лучшую изолированную скорость
стрима~1 примерно с $24.7$ до $48.142$ миллиона кандидатов/с. В Nsight время
изменённого ядра сократилось с $9.678$ до $3.232$ с за 1,202 запуска
($2.99\times$). Что важнее, тот же production-запуск глубины 20 сократился с
$115.334$ до $60.623$ с, то есть ускорился в $1.90\times$. Старый PyTorch-запуск
глубины 100 исчерпал память на глубине 33 через $329.067$ с; этот незавершённый
запуск не считается базой для ускорения.

\paragraph{Kaggle $2\times$T4.}
Два ранга совместно поддерживают один глобальный луч. Для задачи 20 запрос
$B_{\rm req}=82{,}615{,}524$ был выровнен до
$B_{\rm eff}=82{,}837{,}504$, а запуск с целевой окрестностью радиуса 4 решил
задачу за $350.402$ с при 128 шардах и $B_\mu=4096$. При
$B=2^{26}$ увеличение числа шардов с 8 до 64 сократило среднее время насыщенных
глубин 7--9 с $69.365$ до $51.732$ с ($1.34\times$), сохранив точный глобальный
отбор. Лучший изолированный результат стрима~1 равен $40.60$ M кандидатов/с;
он не используется как E2E-скорость.

\paragraph{Ёмкость $8\times$A100 и предел scale-out.}
Сводка exact-scale autotune фиксирует стабильное выполнение глубины 8
для трёх экземпляров Megaminx при запросе
$B_{\rm req}=770{,}883{,}178$ с пиком $39{,}745$ MiB на GPU. Runtime-выравнивание
даёт $B_{\rm eff}\geq B_{\rm req}$, но точная эффективная ширина, сырые rank-логи,
runtime-конфигурация, commit, хеш checkpoint и wall time не сохранились. Поэтому
сводка подтверждает capacity-результат на восьми устройствах и номинальный
насыщенный масштаб не менее $18.501$ миллиарда логических кандидатов-потомков,
но не измеренное время A100. Отдельный H200-запуск ниже даёт измерение скорости
на восьми GPU для другой головоломки. Статические таблицы обмена сейчас
допускают не более 128 рангов, но выполнение за пределами восьми GPU не измерено.

\paragraph{Миллиардная глубина на $8\times$H200.}
На одном экземпляре Vast с восемью H200 (по 143,771 MiB) запуск Cube4,
задача 1000, с 24 генераторами запросил $B=2{,}900{,}000{,}000$ и выровнялся
до $B_{\rm eff}=2{,}900{,}361{,}216$, или $362{,}545{,}152$ удерживаемых
состояний на ранг. Все восемь рангов завершили насыщенную глубину 8 за
$931.266$ с. Это $3.114$ миллиона родителей/с или $74.746$ миллиона
номинальных пар родитель--действие/с end to end, включая оценку, глобальную
редукцию, обмен, отбор и обработку истории. Оценщик --- FP16 piece Transformer
для Cube4 с 24 выходами и 3 383 064 параметрами: четыре блока шириной 256,
восемь голов, feed-forward шириной 1024. Префикс SHA-256 манифеста модели ---
\texttt{03f55e5b}; полный хеш указан в манифесте доказательств. Ревизия
реализации --- \texttt{3b756afe}; использованы 64 шарда на GPU, а занятая
память устройства составляла около 141,820 MiB/GPU. При запросе 2.94 млрд
первый обмен счётчиками NCCL завершился ошибкой; 3.00 млрд не прошли статическую
проверку бюджета памяти. Ни один из этих исходов не классифицируется как CUDA
OOM. Запуск задачи 1000 был намеренно остановлен во время глубины 9, поэтому
время глубины 8 --- измерение завершённого шага, а не полного решения.
Нагрузка и модель Cube4 отличаются от запуска Megaminx на двух T4; эти две
скорости не являются сравнением strong scaling.

\section{Сравнение с CayleyPy и результат на Megaminx}
CayleyPy экспериментально показывает, что широкий луч повышает вероятность
успеха и качество решения, но опубликованные примеры остаются однокарточными.
В контролируемом эксперименте на $S_{20}$ ширина $W$ растёт от $2^{10}$
(решения нет) до $2^{20}$ (идеальная длина 190), а в статье о больших кубах
используется $W=2^{24}$, чтобы один агент помещался на одной
GPU~\cite{soibelman2025cayleypy,chervov2025large}. Измеренная удерживаемая
ширина H200, $2{,}900{,}361{,}216$, в $172.875\times$ шире $2^{24}$, в
$145.02\times$ шире крупнейшего найденного рецензируемого удерживаемого луча
(20 миллионов) и в $116.01\times$ шире публичной реализации HATETRIS с 25
миллионами~\cite{hatetrispublic}. Эти отношения сравнивают удерживаемую
ёмкость, а не скорость или сложность решений на разных графах.

\begin{table}[t]
\centering
\small
\caption{Опубликованная зависимость CayleyPy от ширины луча (таблица 8 из
\cite{soibelman2025cayleypy}).}
\label{tab:cayleypy_width}
\begin{tabular}{@{}rrrrrr@{}}
\toprule
$W$ & $2^{10}$ & $2^{12}$ & $2^{14}$ & $2^{16}$ & $2^{20}$ \\
\midrule
Длина & --- & 252 & 202 & 194 & 190 \\
\bottomrule
\end{tabular}
\end{table}

Дополнительно мы измерили Pilgrim и нашу реализацию поиска на одном
состоянии Cube4, одном checkpoint, одной GPU, одинаковом луче и десяти слоях.
Scorer --- piece Transformer с 3 383 064 параметрами, 96 позициями из шести
классов, $d_{\rm model}=256$, четырьмя слоями, восемью головами, шириной
feed-forward 1024 и 24 Q-выходами; inference использует FP16 и microbatch 384.
Таблица~\ref{tab:cube4_pair} показывает десятый слой, где насыщенный frontier
разворачивает $24B$ пар.

\begin{table}[t]
\centering
\small
\setlength{\tabcolsep}{2pt}
\caption{Согласованное вычисление глубины 10 для Cube4 при $B=4{,}194{,}304$.}
\label{tab:cube4_pair}
\resizebox{\columnwidth}{!}{%
\begin{tabular}{@{}llrrr@{}}
\toprule
GPU & Реализация & с/глубину (глубина 10) & Пар/с & Пик VRAM \\
\midrule
T4 & Pilgrim & 198.153 & 508,007 & 2.90 GiB reserved \\
T4 & Наша & 181.282 & 555,293 & 1,949 MiB used \\
P100 & Pilgrim & 314.487 & 320,087 & 2.90 GiB reserved \\
P100 & Наша & 301.955 & 333,378 & 1,949 MiB used \\
\bottomrule
\end{tabular}%
}
\end{table}

\subsection{Инвентаризация подходов и согласованные семейства моделей на одной T4}
Таблица~\ref{tab:four_t4} является инвентаризацией, а не рейтингом скорости
между строками. AlphaCube и DeepCubeA сохраняют свои опубликованные модели
Cube3. Результаты Megaminx образуют два непересекающихся согласованных семейства:
скалярное \texttt{output\_dim=1} (Pilgrim Searcher, нативный CayleyPy,
MultiGPUBeamSearch) и нативное \texttt{output\_dim=24} (Pilgrim QSearcher,
MultiGPUBeamSearch). Отношение допускается только при совпадении checkpoint,
архитектуры, головы, состояния, генераторов, глубины, ширины, dtype и
оборудования. CayleyPy не принимает 24-выходную голову и остаётся несовместимым
без преобразований. В столбце
скорости приводятся только нативно подсчитанные узлы, делённые на solver wall
time; они никогда не заменяются величиной~$B/t$.

\begin{table*}[t]
\centering
\small
\setlength{\tabcolsep}{3pt}
\caption{Точные архитектуры нейронных сетей в приведённых тестах. Число
параметров относится к обучаемой исходной модели до свёртки batch normalization.}
\label{tab:model_architectures}
\resizebox{\textwidth}{!}{%
\begin{tabular}{@{}lrlp{10.8cm}@{}}
\toprule
Семейство модели & Параметры & Где используется & Последовательность блоков \\
\midrule
Cube3 AlphaCube 0.1.6 \texttt{large} & 118,939,666 & AlphaCube & one-hot 324 $\to$ $8\times$[Linear 4096, ReLU, BN] $\to$ Linear 18 \\
Cube3 DeepCubeA official \texttt{cube3} & 14,663,001 & DeepCubeA & one-hot 324 $\to$ Linear 5000, BN, ReLU $\to$ Linear 1000, BN, ReLU $\to$ $4\times$[Linear 1000, BN, ReLU, Linear 1000, BN, сложение residual, ReLU] $\to$ Linear 1 \\
Megaminx scalar MLP & 34,766,849 & Pilgrim Searcher; CayleyPy; MultiGPU & позиция--класс $120\times120$ $\to$ эквивалент EmbeddingBag Linear 2048, BN, ReLU $\to$ Linear 512, BN, ReLU $\to$ $8\times$[Linear 512, BN, ReLU, Linear 512, BN, сложение residual, ReLU] $\to$ Linear 1 \\
Megaminx QMLP2RB & 23,978,008 & Pilgrim QSearcher; MultiGPU & позиция--класс $120\times120$ $\to$ эквивалент EmbeddingBag Linear 1536, BN, ReLU $\to$ Linear 512, BN, ReLU $\to$ $2\times$[Linear 512, BN, ReLU, Linear 512, BN, сложение residual, ReLU] $\to$ Linear 24 \\
Cube4 Q Transformer & 3,383,064 & только согласованный Cube4 & 96 токенов / 6 классов $\to$ проекция + CLS $\to$ $4\times$[pre-LN, 8-головое attention ($d=256$), residual, pre-LN, FFN 256--1024--256 ReLU, residual] $\to$ LN $\to$ Linear 24 \\
\bottomrule
\end{tabular}%
}
\end{table*}

В согласованных строках Megaminx все перечисленные реализации загружают одни
и те же байты: SHA-256 скалярного checkpoint равен
\texttt{7f5071e6...a539b759}, а checkpoint с 24 выходами ---
\texttt{e7bda332...5d2ca670}. Вторая модель не получена из скалярной. Cube4
Transformer образует отдельное согласованное сравнение и не смешивается с
capacity sweep Megaminx.

\begin{table*}[t]
\centering
\small
\setlength{\tabcolsep}{3pt}
\caption{Инвентаризация нативных измерений на одной видимой T4. Строки с
разными контрактами моделей не сравниваются; максимум измерен.}
\label{tab:four_t4}
\resizebox{\textwidth}{!}{%
\begin{tabular}{@{}lllrrrl@{}}
\toprule
Подход & Нагрузка / нативный масштаб & Макс. complete & Время (с) & с/нативный шаг & Узлы/с & Граница/статус \\
\midrule
AlphaCube 0.1.6 & Cube3 \texttt{large}, beam & 1,572,864 & 461.488 & 27.146/глубину & 42,387 & 1,703,936 host SIGKILL \\
DeepCubeA & Cube3 A*, 55 итераций & н/п & 146.474 & 2.663/итерацию & 41,880 & complete, длина 21 \\
Pilgrim Searcher & Megaminx scalar beam, глубина 8 & 5,898,240 & 1,462.557 & 182.820/глубину & не измерено & 6,291,456 OOM \\
Pilgrim QSearcher & Megaminx output-24 beam, глубина 8 & 20,971,520 & 197.358 & 24.670/глубину & не измерено & 23,068,672 OOM \\
нативный CayleyPy & Megaminx scalar beam, глубина 8 & 20,971,520 & 4,258.602 & 532.325/глубину & не измерено & broad: $2^{24}$ complete, $2^{25}$ OOM \\
MultiGPUBeamSearch & Megaminx scalar beam, глубина 8 & 29,360,128 & 2,817.158 & 352.145/глубину & не измерено & 31,457,280 OOM \\
MultiGPUBeamSearch & Megaminx output-24 beam, глубина 8 & 29,360,128 & 116.904 & 14.613/глубину & не измерено & 33,554,432 OOM \\
\bottomrule
\end{tabular}%
}
\end{table*}

\begin{table}[t]
\centering
\small
\caption{Завершённый sweep AlphaCube на 1xT4.}
\label{tab:alpha_curve}
\begin{tabular}{@{}rrrrr@{}}
\toprule
$B$ & Время (с) & с/глубину & Узлы & Узлы/с \\
\midrule
1,024 & 1.034 & 0.052 & 18,720 & 18,108 \\
4,096 & 1.199 & 0.060 & 73,484 & 61,309 \\
16,384 & 5.014 & 0.251 & 282,380 & 56,321 \\
65,536 & 18.543 & 1.091 & 903,448 & 48,722 \\
131,072 & 37.818 & 2.224 & 1,755,416 & 46,418 \\
262,144 & 79.166 & 4.657 & 3,459,352 & 43,698 \\
524,288 & 170.671 & 10.039 & 6,867,224 & 40,237 \\
786,432 & 228.938 & 13.467 & 10,124,136 & 44,222 \\
1,048,576 & 303.870 & 17.875 & 13,269,864 & 43,670 \\
1,310,720 & 382.158 & 22.480 & 16,415,592 & 42,955 \\
1,572,864 & 461.488 & 27.146 & 19,561,320 & 42,387 \\
\bottomrule
\end{tabular}
\end{table}

\begin{table*}[t]
\centering
\small
\setlength{\tabcolsep}{4pt}
\caption{Завершённый sweep нативного CayleyPy на 1xT4 по степеням двойки.
Память --- пик CUDA-аллокатора PyTorch; последняя строка --- первый отказ
широкой сетки.}
\label{tab:cayleypy_curve}
\begin{tabular}{@{}rrrrrr@{}}
\toprule
$B$ & Время (с) & с/глубину & Allocated (MiB) & Reserved (MiB) & Статус \\
\midrule
1,024 & 0.930 & 0.116 & 85.9 & 92.0 & complete \\
2,048 & 1.237 & 0.155 & 96.3 & 112.0 & complete \\
4,096 & 2.189 & 0.274 & 117.3 & 128.0 & complete \\
8,192 & 3.828 & 0.478 & 158.6 & 230.0 & complete \\
16,384 & 6.776 & 0.847 & 241.5 & 292.0 & complete \\
32,768 & 13.722 & 1.715 & 407.5 & 510.0 & complete \\
65,536 & 31.399 & 3.925 & 739.3 & 858.0 & complete \\
131,072 & 63.664 & 7.958 & 772.2 & 1,180.0 & complete \\
262,144 & 114.551 & 14.319 & 834.8 & 1,168.0 & complete \\
524,288 & 219.991 & 27.499 & 962.1 & 2,156.0 & complete \\
1,048,576 & 430.789 & 53.849 & 1,216.7 & 3,234.0 & complete \\
2,097,152 & 745.341 & 93.168 & 1,726.0 & 2,740.0 & complete \\
4,194,304 & 1,320.000 & 165.000 & 2,745.2 & 4,492.0 & complete \\
8,388,608 & 2,489.356 & 311.169 & 5,155.3 & 6,462.0 & complete \\
16,777,216 & 4,741.581 & 592.698 & 10,233.7 & 14,652.0 & complete \\
33,554,432 & --- & --- & 12,922.2 & 14,484.0 & CUDA OOM \\
\bottomrule
\end{tabular}
\end{table*}

Для AlphaCube с/глубину равно solver wall time, делённому на завершённую
глубину поиска из возвращённого решения (20 для первых трёх строк и 17 для
остальных). Для обеих реализаций Megaminx wall time делится на фиксированные
восемь завершённых глубин. Эти средние включают запуск и ранние ненасыщенные
слои и дополняют, а не заменяют, узлы/с.

Самая широкая скалярная точка Pilgrim, $B=5{,}898{,}240$, завершилась за
1,462.557 с (182.820 с/глубину) с пиком устройства 14,433 MiB;
$B=6{,}291{,}456$ дал CUDA OOM. Для output-24 $B=20{,}971{,}520$ завершился за
197.358 с
(24.670 с/глубину) при 12,333 MiB, тогда как $B=23{,}068{,}672$ дал CUDA OOM.
Измеренные интервалы равны соответственно 5,898,240--6,291,456 и
20,971,520--23,068,672.

Финальный MultiGPU-прогон соавтора и скалярное подтверждение v27 теперь дают
измеренные пары complete/OOM для обеих голов. Скалярный запуск завершился при
29,360,128 и дал CUDA OOM при 31,457,280; нативная серия output-24 завершилась
до 29,360,128 и дала CUDA OOM при 33,554,432.

\begin{table}[t]
\centering
\small
\caption{Финальные скалярные граничные точки MultiGPU на 1xT4.}
\label{tab:multigpu_v11_partial}
\resizebox{\columnwidth}{!}{%
\begin{tabular}{@{}rrrrl@{}}
\toprule
$B$ & Время (с) & с/глубину & Пик устройства (MiB) & Статус \\
\midrule
25,165,824 & 3,027.673 & 378.459 & 12,267 & complete \\
27,262,976 & 3,080.206 & 385.026 & 13,241 & complete \\
29,360,128 & 2,817.158 & 352.145 & 14,215 & complete \\
31,457,280 & --- & --- & 14,909 & CUDA OOM \\
\bottomrule
\end{tabular}
}
\end{table}

Скалярный OOM повторился для каждого профиля row budget от 49,152 до 24.
Общая первичная метрика равна полному wall time, делённому на восемь:
352.145 с/глубину; вспомогательное время финальной глубины равно 1,644.13 с.

\begin{table}[t]
\centering
\small
\caption{Финальный нативный output-24 sweep MultiGPU на 1xT4.}
\label{tab:multigpu_output24}
\resizebox{\columnwidth}{!}{%
\begin{tabular}{@{}rrrrl@{}}
\toprule
$B$ & Время (с) & с/глубину & Пик устройства (MiB) & Статус \\
\midrule
1,048,576 & 8.296 & 1.037 & 773 & complete \\
4,194,304 & 22.735 & 2.842 & 2,253 & complete \\
8,388,608 & 41.091 & 5.136 & 4,159 & complete \\
16,777,216 & 77.460 & 9.683 & 7,983 & complete \\
20,971,520 & 96.124 & 12.015 & 9,893 & complete \\
25,165,824 & 111.130 & 13.891 & 11,799 & complete \\
29,360,128 & 116.904 & 14.613 & 13,717 & complete \\
33,554,432 & 1.563 & --- & 14,909 & CUDA OOM \\
\bottomrule
\end{tabular}
}
\end{table}

OOM output-24 повторялся при снижении runtime-профиля row budget от 2,048 до
1. Поле 1.563 с относится к последней неудачной попытке, а не к завершённой
глубине.

\begin{table}[t]
\centering
\scriptsize
\setlength{\tabcolsep}{3pt}
\caption{Согласованный скалярный Megaminx-sweep на 1xT4. Ни одна реализация
не предоставляет валидный счётчик узлов для этих точек.}
\label{tab:megaminx_curve}
\begin{tabular}{@{}rrrrr@{}}
\toprule
& \multicolumn{2}{c}{MultiGPU} & \multicolumn{2}{c}{CayleyPy} \\
\cmidrule(lr){2-3}\cmidrule(l){4-5}
$B$ & Время (с) & с/глубину & Время (с) & с/глубину \\
\midrule
65,536     & 12.537    & 1.567   & 31.399    & 3.925 \\
262,144    & 39.935    & 4.992   & 114.551   & 14.319 \\
1,048,576  & 152.604   & 19.076  & 430.789   & 53.849 \\
4,194,304  & 476.007   & 59.501  & 1,320.000 & 165.000 \\
8,388,608  & 892.948   & 111.619 & 2,489.356 & 311.169 \\
16,777,216 & 1,728.460 & 216.058 & 4,741.581 & 592.698 \\
20,971,520 & 2,179.990 & 272.499 & 4,258.602 & 532.325 \\
25,165,824 & 2,416.150 & 302.019 & ---       & --- \\
\bottomrule
\end{tabular}
\end{table}

AlphaCube материализует все $18B$ дочерних состояний Cube3 в массивах NumPy
на хосте до сортировки; поэтому первый отказ является исчерпанием памяти хоста,
хотя GPU reserved составляет лишь 4.685~GiB. На самых широких завершённых
точках измеренная скорость стабилизируется около 40--44 тысяч узлов/с.
DeepCubeA не имеет постоянного луча: это weighted A*, а его нативным масштабом
служат search/NN batch 10,000/10,000, поэтому приписывать ему максимальный
фронт было бы некорректно. Завершённый T4-запуск выполнил 55 поисковых
итераций за 146.474 с, то есть 2.663 с на одну нативную A*-итерацию; это не
измерение beam-глубины. За запуск обработано 6,134,352 узла со скоростью
41,880 узлов/с и получено решение длины 21.

Для согласованной скалярной пары Megaminx CayleyPy commit \texttt{5b2b53f}
использует iterated search с \texttt{history\_depth=0}. Широкий протокол
завершил каждую степень двойки от $2^{10}$ до $2^{24}$ и получил CUDA OOM на
$2^{25}$; оценка по тренду памяти совпала с $2^{24}$, поэтому лишняя
подтверждающая точка не запускалась. Более раннее целевое измерение с тем же
контрактом сохраняет более узкую пару 20,971,520 complete против 21,037,056
CUDA OOM, но не выдаётся за часть широкой сетки. Внутренний путь
расширения обошёл публичный hook счётчика, поэтому валидны только wall time и
память. MultiGPU также сохранил end-to-end wall, но не валидный счётчик узлов.
Финальный скалярный прогон измерил 29,360,128 complete при 14,215 MiB против
31,457,280 CUDA OOM при 14,909 MiB. Согласованный прогон output-24 измерил
29,360,128 complete против 33,554,432 CUDA OOM.
Отказы устаревшего статического профиля и отменённые v10/v11 исключены. В
последней согласованной
точке широкой сетки CayleyPy занял 4,741.581 с, а исторический MultiGPU-запуск
1,728.460 с при $B=16{,}777{,}216$ (wall ниже в 2.74 раза). На более ранней
целевой точке $B=20{,}971{,}520$ wall MultiGPU ниже в 1.95 раза; это отношения
согласованного end-to-end wall, а не восстановленные отношения nodes/s.

Эти скалярные отношения используют один идентичный контракт скалярной модели и
не сравниваются с более быстрой строкой Pilgrim QSearcher. 24-выходная модель
оценивает все 24 действия для каждого родителя и имеет другую архитектуру;
поэтому разница scalar--Q является разницей контракта модели, а не ускорением
алгоритма. Отдельное сравнение Cube4 использует идентичный 24-выходный
Q-checkpoint для Pilgrim и MultiGPUBeamSearch. Значение output-24 для CayleyPy
не выводится.

Таким образом, в согласованной
точке с Pilgrim наша пропускная способность на
десятом слое выше в $1.093\times$ на T4 и в $1.042\times$ на P100. Резерв
аллокатора PyTorch и полное использование устройства по \texttt{nvidia-smi}
являются разными метриками памяти, поэтому таблица не утверждает отношение
расхода памяти; она устанавливает, что обе реализации завершают луч 4M.

Для Megaminx-superflip найденная публичная хронология улучшается от 55 через 54
и 53 до 51 хода; следовательно, непосредственно предшествовавшее публичное
значение было 53, а не 55~\cite{megaminxsuperflip53,megaminxsuperflip51}.
Путь из 51 шага использует ту же метрику 24 знаковых поворотов одной грани на
72 градуса и имеет публичный wall time $22{,}600.2$ с. Независимый CPU replay
сохранённого submission достигает точного центрального состояния с нулём различий;
SHA-256 пути и файла записаны в сопровождающем аудите доказательств.
В цитируемой публичной хронологии нет более короткого пути в той же метрике, но
формального реестра и доказательства оптимальности нет. Поэтому мы называем его
кратчайшим найденным публичным решением в той же метрике, а не безусловным
мировым рекордом. Отсутствующий манифест этого запуска отделён от capacity-сводки
глубины 8 для запроса $770{,}883{,}178$ состояний.

\section{Код и артефакты}
Исходники статьи, индекс доказательств и машиночитаемые таблицы общедоступны по
адресу \url{https://github.com/TryDotAtwo/MultiGPUBeamSearchPaper}. Базовый
релиз v1.0.0 содержит согласованные checkpoints и предыдущие PDF-версии:
\url{https://github.com/TryDotAtwo/MultiGPUBeamSearchPaper/releases/tag/v1.0.0}.
Исходники solver и его готовый к запуску на GPU Megaminx-релиз общедоступны по
адресам \url{https://github.com/TryDotAtwo/MultiGPUBeamSearch} и
\url{https://github.com/TryDotAtwo/MultiGPUBeamSearch/releases/tag/megaminx-native-cluster-latest}.
Публичный самодостаточный одно-T4 датасет для воспроизведения доступен по адресу
\url{https://www.kaggle.com/datasets/trydotatwo/gpu-beam-search-paper-1xt4-reviewer-notebooks}.
Четыре импортируемых исходника ноутбуков для широких sweep и подтверждений
границ Pilgrim и MultiGPUBeamSearch, а также runners CayleyPy, AlphaCube и
DeepCubeA, версионируются в репозитории статьи. Они используют публичный датасет
через обычный workflow Kaggle Import Notebook и Save Version и не зависят от
видимости сохранённой страницы kernel владельца.
Запись для восьми A100 даёт доказательства ёмкости и памяти, но не wall time или
точную эффективную ширину; запись для двух T4 отдельно даёт измеренную
пропускную способность. Запуск Cube4 на восьми H200 дополнительно даёт
завершённую миллиардную глубину с полными rank-логами, телеметрией устройств
и манифестом доказательств по адресу
\url{https://github.com/TryDotAtwo/MultiGPUBeamSearchPaper/tree/main/artifacts/results/h200_8x_cube4}.
Манифест фиксирует реализацию, модель и прерванную следующую глубину.
Отношение ускорения между этими разными нагрузками не выводится.

\section{Заключение}
Эта работа устанавливает один GPU-резидентный распределённый beam search, а не
набор локальных суррогатов. Для фиксированной раскладки его потоковый контур
expand--deduplicate--select, без семантических локальных ограничений шардов,
математически эквивалентен монолитному глобальному top-$B$ по редуцированной
оценке ключа при той же квантизации, порядке разрешения равных оценок и
заявленной 128-битной модели идентичности. Результат Cube4 на восьми H200
устанавливает физически распределённый удерживаемый frontier из
$2{,}900{,}361{,}216$ состояний и измеряет одну завершённую насыщенную глубину
со скоростью $74.746$ миллиона номинальных пар родитель--действие/с end to end.
Отдельный результат Megaminx на двух T4 измеряет $30.274$ миллиона логических
потомков/с; разные нагрузки не доказывают strong scaling. Сводка ёмкости для
восьми A100 не содержит сохранённого времени. Проект также
предоставляет replay-valid путь Megaminx-superflip длиной 51 ход --- кратчайшее
публичное решение в той же метрике, установленное сопровождающими
доказательствами. Это первая опубликованная высокопроизводительная multi-GPU
система beam search с одним физически шардированным глобальным frontier и
заявленным контрактом точности при масштабе удерживаемого frontier в миллиарды
состояний.

\begingroup
\small
\bibliographystyle{abbrv}
\bibliography{references}
\endgroup
\end{document}
